\scanpage{268}
\mspage{9}

Ainsi les deux catégories ($\unsure{\Psi|\Phi}_1$ et $\Phi_1$) ci-dessus sont symétriques, mais opposées l'une de l'autre sans être isomorphes. D'autre part \deleted{$\mathbb{\Delta}$} les $\mathbb{\Delta}_i$ ne sont pas symétriques, mais \unsure{chacuns|chacun|chacune} \unsure{isomorphes|isomorphe} à la catégorie opposée.

Si $I$ est isomorphe à la catégorie opposée $I^{\unsure{\circ|\mathrm{o}}}$, et si $\Theta$ est un tel isomorphisme (\unsure{néc.} \deleted{\unsure{pr}}unique), alors \deleted{\unsure{il sera}} \deleted{il est immédiat} il est immédiat que $\Theta$ renverse l'orientation (sinon ce serait l'identité sur les sommets, et les flèches de transition seraient à la fois directes et indirectes), \deleted{de} donc $\Theta(0)=n$, $\Theta(1)=n-1$, et on voit que $u_1$ et $u_n$ sont de même signe. Mais si $I$ est symétrique, $u_1$ et $u_n$ sont de signes opposés. Donc $I$ ne peut être à la fois symétrique, et isomorphe à \deleted{\unsure{$I^\circ$}} l'intervalle opposé.

Pour toute catégorie $\mathbb{\Delta}_i$, celle-ci étant un intervalle, il y a \inserted{(si $i>0$)} \uline{deux orientations}. \deleted{L'orientation} En tant qu'intervalles orientés, il y a $\mathbb{\Delta}_i^+$ (avec l'\deleted{relation}orientation ordinaire, donnée par l'ordre ordinaire sur $\operatorname{Ob}\mathbb{\Delta}_i=[0,i]$), et $\mathbb{\Delta}_i^-$ avec l'orientation opposée.

\deleted{\unsure{Ce|On}}

\deleted{\uline{Proposition} Soit $I$ une \deleted{\unsure{inter}} catégorie intervalle orientée, (de longueur $n\geq1$) donc can. isom. à un intervalle standart. Alors il y a une suite unique de}

\[\deleted{\mathbb{\Delta}_{l_1},\ldots,\mathbb{\Delta}_{l_r}}\]

\deleted{de simplexes types, \deleted{de long} de dimensions $l_1,\ldots l_r>0$, \deleted{chacun orienté} et une suite unique d'orientations de ceux-ci, de telle façon que}

% archive sheet number: 267
